\section{Closing the latency gap: \texttt{fast\_lag} from stage calibration to eight-A100 wall-clock speedup}
\label{sec:systems}

At the algorithmic level, lagged output normalization should add very little
wall-clock cost. Vanilla input-side \taxid{MuonT} (\muont{}) computes three
width-512 and one width-2048 dense moments per layer plus three right matrix
products per update, yet runs only $+\sysMuontVsMuon\%$ ($+21.2$\,ms/step in
Wave~0; $+\sysExecMuontVsMuonPct\%$ in \texttt{wave\_sym\_fastlag}) slower than
plain \taxid{Muon} because forward row capture uses three fused Triton launches
and input root refinement overlaps the backward pass. By contrast,
\taxid{WT-MLP} (\mlpdiag{}) adds five width-512 dense moments and \emph{zero}
width-2048 dense moments per layer --- less nominal $O(d^2 + d^3)$ moment and
eigendecomposition work than \taxid{MuonT}'s own input side --- and one-update
lag (\taxid{WT-MLP-Lag}, \taxid{WT-Sym-Lag}; Figure~\ref{fig:method_arch}b)
makes roughly $900$\,ms of the next step's forward and backward pass available
to hide output statistics and root work.

Why, then, did Wave~4's \taxid{WT-MLP-Lag} run $+\sysLegacyMlpVsMuon\%$ slower
than \taxid{Muon} ($+18.0\%$ over \taxid{MuonT}), and Wave~3's
\taxid{DT-All-Lag} (\denselag{}) run $+\sysLegacyAllVsMuon\%$ slower?
Stage-level profiling alongside Wave~10's matched null (\taxid{Null-OD} /
\nullod{}, which still cost $+\sysWaveTenNullVsMuon\%$ over \taxid{Muon} with
\emph{zero} dense output roots) localized the gap to implementation bottlenecks
rather than intrinsic arithmetic cost. We detail the three-pillar redesign
(\texttt{fast\_lag}) below, followed by its full-budget execution across four
uncontended eight-A100 nodes in \texttt{wave\_sym\_fastlag}.

\subsection{Root-cause decomposition and three-pillar redesign}
\label{sec:systems-pillars}

\paragraph{Pillar 1: True asynchronous lag and opportunistic mid-step fallback resolution.}
Two synchronization barriers in the initial runner prevented one-update lag
from overlapping with the next step:
\begin{enumerate}
\item \texttt{train.py} executed a device-wide \texttt{torch.cuda.synchronize()}
  at the end of every update to record \texttt{step\_seconds}, stalling the
  host until the newly launched output side stream completed inside the
  \emph{current} update.
\item Even more expensively, \texttt{OutputToner.begin()} only enqueued
  candidate graph roots; flagged root-repair iterations and exact
  \texttt{cuSOLVER eigh} fallbacks ($\approx 4$--$6$ per update globally) were
  deferred until \texttt{\_collect\_lagged()} invoked \texttt{\_finish\_roots()}
  right before \texttt{optimizer.step()} on update $t+1$, blocking the main
  stream synchronously ($+\sysRootsLegacyMlpMs$\,ms in \texttt{mlpdiag\_lag}
  and $+\sysRootsLegacyAllMs$\,ms in \texttt{all\_lag}).
\end{enumerate}
In \texttt{fast\_lag}, end-of-step timing synchronizes only
\texttt{torch.cuda.current\_stream()}, while full side-stream synchronization
occurs only at checkpoint serialization and shutdown. To prevent update $t+1$'s
first backward pass from overwriting update $t$'s sampled cotangent buffers
without doubling memory, a lightweight CUDA event \texttt{ev\_rows\_consumed} is
recorded on the priority-0 output side stream immediately after
\texttt{syrk\_list\_packed} and local diagonal row-sums finish reading the
buffers ($\approx 1.4$\,ms); \texttt{select\_microbatch(0)} on update $t+1$
waits on that event at $0.0$\,ms stall. Once microbatch~0's $\approx 235$\,ms
forward--backward pass is queued on the main stream,
\texttt{select\_microbatch($k\ge 1$)} opportunistically polls readiness and
resolves all repairs, exact \texttt{eigh} fallbacks, unit-RMS normalizations,
and FP16 factor splits on the background side stream roughly $700$\,ms before
\texttt{step()} consumes them.

\paragraph{Pillar 2: Per-site block-level Triton capture and symmetric inline diagonals.}
Wave~10's null arm revealed a $+\sysCapLegacyMs$\,ms/step floor even with no
dense output roots, caused by wrapping all 224 linear layers in
\texttt{RowObservedLinear} ($896$ opaque custom-op boundaries per update inside
\texttt{torch.compile}'s AOTAutograd backward graph and $2{,}688$ eager gather
kernels writing $1.8$\,GiB of BF16 rows). We replace all 224 wrappers with four
view-transparent block-level autograd hooks per transformer block. When
$(\text{gate},\text{up})$ use diagonal moments, an atomic-free 2D-tiled Triton
kernel (\texttt{BLOCK\_R=16, BLOCK\_C=64}) accumulates sampled FP32 row squares
directly in registers without allocating the 64 $[4576, 2048]$ BF16 buffers ---
cutting backward capture overhead $\sysCapSpeedup\times$ (from
$+\sysCapLegacyMs$\,ms to $+\sysCapFastMlpMs$\,ms) and saving
$\sysMemSavedGib$\,GiB in isolated 32-layer profiles ($\sysExecMlpLagSavedWaveFourGib$\,GiB
in full eight-rank DDP). Under eight-rank \texttt{DDP(gradient\_as\_bucket\_view=True)}
with \texttt{torch.compile}, DDP's 25\,MB gradient buckets split \texttt{model.body}
into $\approx 2$-layer compiled backward partitions; allocating independent
per-site FP32 partial tensors (\texttt{InlineDiagonalAccumulator}, $128$\,KB per
block, $4$\,MiB total, cleared in one \texttt{torch.\_foreach\_zero\_} call)
prevents cross-bucket saved-tensor \texttt{\_version} increments while
preserving bitwise-exact reductions. In \taxid{WT-Sym} and \taxid{WT-Sym-Lag},
the same deterministic two-stage Triton \texttt{reduce\_row\_squares} +
$262$\,KB \texttt{all\_reduce} + \texttt{diagonal\_roots} pipeline runs
asynchronously on \texttt{\_input\_side\_stream} for width-2048 \texttt{in:down}
during \texttt{begin()}, deleting \texttt{self.input.groups["f"]} and
eliminating all $2048\times 2048$ input \texttt{syrk}, $268.6$\,MB
\texttt{reduce\_scatter}, cached eigenbases, and host \texttt{eigh} jobs.

\paragraph{Pillar 3: Compiled Ozaki FP16-split Tensor-Core update stack.}
Because deterministic training disables TF32 (\texttt{allow\_tf32=False}),
eager FP32 \texttt{torch.bmm} products executed on $19.5$\,TFLOPS SIMT cores and
synced \texttt{.tolist()} host finite checks every update. Compiling
\texttt{block\_update\_stack} (supporting both 2D dense and 1D diagonal right
and left factors in \texttt{\_compiled\_block\_in} / \texttt{\_compiled\_block\_out})
with upstream's FP16-split Ozaki scheme (\texttt{PR.\_EXACT}) moves all dense
factor products onto $312$\,TFLOPS Tensor Cores with full FP32 accuracy, fuses
post-$\Nten$ direction and decay into two right products instead of three, and
masks nonfinite updates on device, shrinking owned-matrix update overhead from
$+\sysUpdLegacyMlpMs$\,ms to $+\sysUpdFastMlpMs$\,ms.

\begin{table}[t]
\centering
\caption{Eight-A100 steady-step latency decomposition and peak Rank-0 allocation
across legacy DDP waves, calibrated single-process SM80 ABBA stage profiles, and
executed uncontended eight-A100 DDP full-budget runs (\texttt{wave\_sym\_fastlag},
steps $2{,}049$--$4{,}096$ median). Under full eight-rank DDP, background
\texttt{reduce\_scatter} and mid-step \texttt{eigh} fallbacks share NVLink and
SMs with DDP's 25\,MB bucket \texttt{all\_reduce} collectives, adding
$\approx 2\%$--$4\%$ over single-process stage calibration; removing
\texttt{in:down}'s $268.6$\,MB \texttt{reduce\_scatter} in \taxid{WT-Sym-Lag}
reclaims $\sysExecSymLagSavedMlpLagMs$\,ms/step and $\sysExecSymLagSavedMlpLagGib$\,GiB.}
\label{tab:systems}
\scriptsize\setlength{\tabcolsep}{3.8pt}
\resizebox{\textwidth}{!}{% Generated by scripts/analyze_systems.py. Do not edit.
\begin{tabular}{@{}llrrrrrrr@{}}
\toprule
Taxonomic ID & Runner / regime & Capture & Update & Async / DDP & Step (s) & vs.\ Muon & vs.\ MuonT & Peak GiB \\
\midrule
\taxid{Muon} & Wave 0 (\texttt{Ref}) & $-1.8$\,ms & --- & $0.0$\,ms & $1.074$ & $0.0$\% & $-1.9$\% & $18.97$ \\
\taxid{MuonT} & Wave 0 / \texttt{sym\_fastlag} & $0.0$\,ms & $0.0$\,ms & $0.0$\,ms & $1.095$ / $1.099$ & $+2.0$\% / $+2.3$\% & $0.0$\% & $20.53$ \\
\addlinespace[2.5pt]
\multicolumn{9}{@{}l}{\textcolor{famAbl}{\pagesfsb Legacy 8$\times$A100 DDP Runner (Waves 3, 4, 10)}} \\
\taxid{DT-All-Lag} & Legacy (Wave 3) & $+92.9$\,ms & $+19.8$\,ms & $+230$\,ms & $1.438$ & $+33.9$\% & $+25.6$\% & $24.49$ \\
\taxid{WT-MLP-Lag} & Legacy (Wave 4) & $+92.9$\,ms & $+15.0$\,ms & $+89.1$\,ms & $1.292$ & $+20.3$\% & $+18.0$\% & $22.68$ \\
\taxid{Null-OD} & Legacy (Wave 10) & $+96.0$\,ms & $0.0$\,ms & $0.0$\,ms & $1.195$ & $+11.3$\% & $+9.1$\% & $22.39$ \\
\addlinespace[2.5pt]
\multicolumn{9}{@{}l}{\textcolor{famDT}{\pagesfsb Calibrated Single-Process SM80 Stage Profile (\texttt{fast\_lag})}} \\
\taxid{DT-All-Lag} & \texttt{fast\_lag} (stage cal.) & $+10.4$\,ms & $+2.3$\,ms & $+21.0$\,ms & $1.129$ & $+4.8\text{--}5.4$\% & $+2.8\text{--}3.4$\% & $24.55$ \\
\taxid{WT-MLP-Lag} & \texttt{fast\_lag} (stage cal.) & $+11.1$\,ms & $+1.7$\,ms & $+6.6$\,ms & $1.115$ & $+3.6\text{--}4.1$\% & $+1.6\text{--}2.1$\% & $21.58$ \\
\addlinespace[2.5pt]
\multicolumn{9}{@{}l}{\textcolor{famWT}{\pagesfsb Executed Uncontended 8$\times$A100 DDP Full-Budget Wave (\texttt{wave\_sym\_fastlag})}} \\
\taxid{WT-Sym} & Synchronous ($\tau=0$) & $+11.1$\,ms & $+1.7$\,ms & $+123.0$\,ms & $1.234$ & $+14.9$\% & $+12.4$\% & $20.73$ \\
\rowcolor{pagepanel} \textbf{\taxid{WT-MLP-Lag}} & \textbf{\texttt{fast\_lag} 8$\times$A100} & $+11.1$\,ms & $+1.7$\,ms & $+50.4$\,ms & $\mathbf{1.162}$ & $\mathbf{+8.2}$\textbf{\%} & $\mathbf{+5.8}$\textbf{\%} & $\mathbf{21.59}$ \\
\rowcolor{pagepanel} \textbf{\taxid{WT-Sym-Lag}} & \textbf{\texttt{fast\_lag} 8$\times$A100} & $+11.1$\,ms & $+1.7$\,ms & $\mathbf{+29.4}$\,\textbf{ms} & $\mathbf{1.141}$ & $\mathbf{+6.2}$\textbf{\%} & $\mathbf{+3.8}$\textbf{\%} & $\mathbf{20.81}$ \\
\bottomrule
\end{tabular}
}
\end{table}

\begin{figure}[t]
\centering
\resizebox{\textwidth}{!}{%
\begin{tikzpicture}
\begin{axis}[
  paperplot,
  width=0.86\textwidth, height=6.8cm,
  xbar stacked, bar width=9.0pt,
  title={8-A100 steady-step latency breakdown above vanilla \taxid{Muon} ($1{,}074$\,ms base)},
  xlabel={Extra step latency above \taxid{Muon} (ms/update; $1{,}074$\,ms base)},
  xmin=0, xmax=535,
  ytick={0,1,2,3,4,5,6,7},
  yticklabels={
    {\taxid{MuonT} ($+2.0\text{--}2.3\%$)},
    {\taxid{DT-All-Lag} legacy ($+33.9\%$)},
    {\taxid{WT-MLP-Lag} legacy ($+20.3\%$)},
    {\taxid{WT-Sym} no-lag 8$\times$A100 ($+14.9\%$)},
    {\taxid{Null-OD} Wave 10 ($+11.3\%$)},
    {\taxid{DT-All-Lag} \texttt{fast} cal.\ ($+5.1\%$)},
    {\textbf{\taxid{WT-MLP-Lag}} \texttt{fast} 8$\times$A100 ($\mathbf{+8.2\%}$)},
    {\textbf{\taxid{WT-Sym-Lag}} \texttt{fast} 8$\times$A100 ($\mathbf{+6.2\%}$)}
  },
  legend columns=4,
  legend style={at={(0.5,-0.23)}, anchor=north, column sep=6pt},
]
\addplot[fill=famRef!35, draw=famRef!70] table[y=idx, x=input_base_ms] {figs/systems_bars.dat};
\addlegendentry{Input \taxid{MuonT} base}
\addplot[fill=famNull!75, draw=famNull!95] table[y=idx, x=capture_ms] {figs/systems_bars.dat};
\addlegendentry{Bwd capture}
\addplot[fill=famOD!75, draw=famOD!95] table[y=idx, x=update_ms] {figs/systems_bars.dat};
\addlegendentry{Ozaki GEMMs}
\addplot[fill=famWT!78, draw=famWT!95] table[y=idx, x=roots_ms] {figs/systems_bars.dat};
\addlegendentry{Exposed async roots / DDP overlap}
\node[font=\tiny, text=famRef, anchor=west] at (axis cs:24,0) {$+21.2$\,ms ($+24.7$\,ms in \texttt{sym\_fastlag})};
\node[font=\tiny, text=famAbl!90!black, anchor=west] at (axis cs:366,1) {$+364$\,ms ($+25.6\%$ vs.\ \taxid{MuonT})};
\node[font=\tiny, text=famAbl!90!black, anchor=west] at (axis cs:221,2) {$+218$\,ms ($+18.0\%$ vs.\ \taxid{MuonT})};
\node[font=\tiny, text=famWT!90!black, anchor=west] at (axis cs:163,3) {$+161$\,ms ($+\sysExecSymVsMuontPct\%$ vs.\ \taxid{MuonT}, $\tau=0$)};
\node[font=\tiny, text=famNull!90!black, anchor=west] at (axis cs:120,4) {$+121$\,ms ($+9.1\%$ vs.\ \taxid{MuonT})};
\node[font=\tiny, text=famDT!95!black, anchor=west] at (axis cs:58,5) {$+55$\,ms (stage cal.: $+3.1\%$ vs.\ \taxid{MuonT})};
\node[font=\tiny, text=famWT, anchor=west] at (axis cs:91,6) {$\mathbf{+88}$\,\textbf{ms} ($\mathbf{+\sysExecMlpLagVsMuontPct\%}$ vs.\ \taxid{MuonT}; cal.\ $+1.8\%$)};
\node[font=\tiny, text=famWT, anchor=west] at (axis cs:70,7) {$\mathbf{+67}$\,\textbf{ms} ($\mathbf{+\sysExecSymLagVsMuontPct\%}$ vs.\ \taxid{MuonT}; \textbf{fastest to $\text{C4}\le 3.52$})};
\end{axis}
\end{tikzpicture}%
}
\caption{Per-update latency above vanilla \taxid{Muon} ($1.074$\,s) from legacy
Waves~3--10 through single-process SM80 stage calibration to executed
uncontended eight-A100 DDP runs (\texttt{wave\_sym\_fastlag}). On full eight-rank
DDP, \texttt{fast\_lag} cuts \taxid{WT-MLP-Lag} step overhead over \taxid{MuonT}
by more than $3\times$ ($+17.6\% \to +\sysExecMlpLagVsMuontPctTwo\%$, saving
$\sysExecMlpLagSavedWaveFourMs$\,ms/step and $\sysExecMlpLagSavedWaveFourS$\,s
end-to-end) and cuts \taxid{WT-Sym-Lag} from $+\sysExecSymVsMuontPctTwo\%$
($\tau=0$) to $+\sysExecSymLagVsMuontPctTwo\%$ (saving
$\sysExecSymLagSavedNolagMs$\,ms/step and $\sysExecSymLagSavedNolagS$\,s
end-to-end).}
\label{fig:systems}
\end{figure}

\subsection{Parity verification, eight-A100 DDP execution, and wall-clock crossings}
\label{sec:systems-parity}

All optimizations and symmetric-MLP extensions pass every correctness gate on
SM80 A100 hardware: all \sysCpuUnitTests{} CPU unit tests, all 11 two-rank NCCL
\texttt{experiment.gpu\_checks} gates plus a compiled 32-layer DDP multi-step
verification, pre- and post-$\Nten$ FP64 reference checks ($\max$ relative error
$4.16\times 10^{-6}$), \texttt{tpu}-mode update parity against upstream
\texttt{FastTonerCore} ($\max$ relative error $1.02\times 10^{-7}$), and
24-update uninterrupted-versus-step-22-resumed continuation past the 20-step
exact-root warmup and CUDA-graph transition with \emph{zero} bitwise error
($0.00$ across all 96 tensors in \mlpdiaglag{}, 99 in \denselag{},
\sysContSymNolagTensors{} in \symmlpdiag{}, and \sysContSymLagTensors{} in
\symmlpdiaglag{}).

Tables~\ref{tab:systems} and~\ref{tab:crossings} and Figure~\ref{fig:systems}
report the executed four-node uncontended eight-A100 wave
(\texttt{wave\_sym\_fastlag}). Three systems findings stand out:
\begin{enumerate}
\item \textbf{More than $3\times$ full-run DDP overhead reduction on
  \taxid{WT-MLP-Lag}:} Relative to same-wave \taxid{MuonT} ($\sysExecMuontStepS$\,s/step,
  $\sysExecMuontEtoeS$\,s end-to-end, $\sysExecMuontGib$\,GiB), \texttt{fast\_lag}
  \taxid{WT-MLP-Lag} runs at $\sysExecMlpLagStepS$\,s/step
  ($+\sysExecMlpLagVsMuontPctTwo\%$ step, $+\sysExecMlpLagEtoePctTwo\%$ end-to-end
  at $\sysExecMlpLagEtoeS$\,s; $+\sysExecMlpLagVsMuonPctTwo\%$ step and
  $+\sysExecMlpLagEtoeMuonPctTwo\%$ end-to-end over Wave-0 \taxid{Muon}) and
  $\sysExecMlpLagGib$\,GiB peak allocation --- cutting Wave~4's legacy overhead
  ($+17.58\%$ vs.\ same-wave \taxid{MuonT}) by $\sysExecMlpLagSavedWaveFourMs$\,ms/step,
  $\sysExecMlpLagSavedWaveFourS$\,s ($11.4$\,min) end-to-end, and
  $\sysExecMlpLagSavedWaveFourGib$\,GiB of VRAM, closely matching the isolated
  $\sysMemFastMlpGib$\,GiB memory projection.
\item \textbf{Calibrated single-process versus full eight-rank DDP overlap:}
  The $\approx 3.9$-percentage-point difference between single-process stage
  calibration ($+\sysFastMlpVsMuontMid\%$ vs.\ \taxid{MuonT}) and full eight-rank
  DDP execution ($+\sysExecMlpLagVsMuontPctTwo\%$) reflects multi-rank NVLink and
  SM contention between the priority-0 output side stream's \texttt{reduce\_scatter}
  / mid-step \texttt{cuSOLVER eigh} fallbacks ($\approx 20{,}000$--$22{,}000$ per
  rank) and DDP's 25\,MB bucket \texttt{all\_reduce} collectives during backward
  microbatches $1..3$. Eliminating dense $2048\times 2048$ \texttt{in:down}
  \texttt{reduce\_scatter} in \taxid{WT-Sym-Lag} removes $\sysExecSymLagSavedMlpLagMs$\,ms/step
  of that contention, bringing full-run eight-A100 overhead down to
  $\sysExecSymLagStepS$\,s/step ($+\sysExecSymLagVsMuontPctTwo\%$ step and
  $+\sysExecSymLagEtoePctTwo\%$ end-to-end over same-wave \taxid{MuonT};
  $+\sysExecSymLagVsMuonPctTwo\%$ step and $+\sysExecSymLagEtoeMuonPctTwo\%$
  end-to-end over \taxid{Muon}) at $\sysExecSymLagGib$\,GiB --- saving
  $\sysExecSymLagSavedNolagMs$\,ms/step and $\sysExecSymLagSavedNolagS$\,s over
  synchronous no-lag \taxid{WT-Sym} ($\sysExecSymStepS$\,s/step,
  $+\sysExecSymVsMuontPctTwo\%$).
\item \textbf{Uncensored launch-anchored wall-clock win at $\text{C4}\le 3.52$:}
  Table~\ref{tab:crossings} reports launch-anchored elapsed seconds (covering
  compilation, training, and periodic evaluation) across all three C4 milestones.
  At $\text{C4}\le 3.59$, \taxid{WT-Sym-Lag} ($\sysCrossSymLagThreeFiveNineS$\,s)
  already beats vanilla \taxid{Muon} ($4{,}880.0$\,s). At $\text{C4}\le 3.52$,
  \taxid{WT-Sym-Lag} crosses at update $4{,}501$ in
  $\sysCrossSymLagThreeFiveTwoS$\,s --- fastest overall ahead of \taxid{Muon}
  ($\sysCrossMuonWaveZeroThreeFiveTwoS$\,s), same-wave \taxid{MuonT}
  ($\sysCrossMuontSymThreeFiveTwoS$\,s), and Wave-0 \taxid{MuonT}
  ($\sysCrossMuontWaveZeroThreeFiveTwoS$\,s) --- while \texttt{fast\_lag}
  \taxid{WT-MLP-Lag} ($\sysCrossMlpLagFastThreeFiveTwoS$\,s) also edges Wave-0
  \taxid{MuonT}. At $\text{C4}\le 3.46$, \taxid{Muon} is censored at budget
  exhaustion ($\waveZeroMuonCfourFinalAbs{}$), whereas all \texttt{WT} arms cross
  at update $5{,}001$ and finish $0.0021$--$0.0029$ nats below same-wave
  \taxid{MuonT}.
\end{enumerate}

\begin{table}[t]
\centering
\caption{Uncensored launch-anchored elapsed-time crossings ($\text{C4}\le 3.59,
3.52, 3.46$) and full-budget endpoints across eight-A100 runs. Every \texttt{WT}
arm reaches $\text{C4}\le 3.52$ in $250$ fewer updates ($2.36\text{B}$ versus
$2.49\text{B}$ tokens); on \texttt{fast\_lag}, \taxid{WT-Sym-Lag} converts that
token saving into the fastest absolute wall-clock crossing
($\sysCrossSymLagThreeFiveTwoS$\,s), while vanilla \taxid{Muon} is censored at
$\text{C4}\le 3.46$.}
\label{tab:crossings}
\scriptsize\setlength{\tabcolsep}{4.5pt}
\resizebox{\textwidth}{!}{% Generated by scripts/analyze_systems.py. Do not edit.
\begin{tabular}{@{}lrrrrr@{}}
\toprule
Arm (wave / runner) & $\text{C4}\le 3.59$ & $\text{C4}\le 3.52$ & $\text{C4}\le 3.46$ & Final C4 & Total E2E \\
\midrule
\taxid{Muon} {\scriptsize(\texttt{wave0})} & $4{,}880.0$\,s {\scriptsize(step 4{,}251)} & $5{,}467.0$\,s {\scriptsize(step 4{,}751)} & \textcolor{famAbl}{\textbf{Censored}} {\scriptsize($3.4618$)} & $3.4618$ & $5{,}889.7$\,s \\
\taxid{MuonT} {\scriptsize(\texttt{wave0})} & $4{,}721.3$\,s {\scriptsize(step 4{,}001)} & $5{,}566.3$\,s {\scriptsize(step 4{,}751)} & $5{,}848.2$\,s {\scriptsize(step 5{,}001)} & $3.4527$ & $5{,}984.6$\,s \\
\taxid{MuonT} {\scriptsize(\texttt{wave\_sym\_fastlag})} & $4{,}619.6$\,s {\scriptsize(step 4{,}001)} & $5{,}468.4$\,s {\scriptsize(step 4{,}751)} & $5{,}754.4$\,s {\scriptsize(step 5{,}001)} & $3.4511$ & $5{,}891.3$\,s \\
\rowcolor{pagepanel} \textbf{\taxid{WT-Sym-Lag}} {\scriptsize(\texttt{fast\_lag})} & $4{,}869.6$\,s {\scriptsize(step 4{,}001)} & $\mathbf{5{,}458.9}$\,\textbf{s} {\scriptsize(step 4{,}501)} & $6{,}040.3$\,s {\scriptsize(step 5{,}001)} & $\mathbf{3.4482}$ & $6{,}179.3$\,s \\
\rowcolor{pagepanel} \textbf{\taxid{WT-MLP-Lag}} {\scriptsize(\texttt{fast\_lag})} & $4{,}966.4$\,s {\scriptsize(step 4{,}001)} & $5{,}564.6$\,s {\scriptsize(step 4{,}501)} & $6{,}157.2$\,s {\scriptsize(step 5{,}001)} & $3.4491$ & $6{,}300.3$\,s \\
\taxid{WT-Sym} {\scriptsize(no-lag, $\tau=0$)} & $5{,}187.7$\,s {\scriptsize(step 4{,}001)} & $5{,}822.6$\,s {\scriptsize(step 4{,}501)} & $6{,}446.1$\,s {\scriptsize(step 5{,}001)} & $\mathbf{3.4482}$ & $6{,}593.9$\,s \\
\taxid{WT-MLP-Lag} {\scriptsize(legacy Wave 4)} & $5{,}500.6$\,s {\scriptsize(step 4{,}001)} & $6{,}164.4$\,s {\scriptsize(step 4{,}501)} & $6{,}824.2$\,s {\scriptsize(step 5{,}001)} & $\mathbf{3.4480}$ & $6{,}985.7$\,s \\
\taxid{WT-MLP} {\scriptsize(legacy Wave 2)} & $5{,}624.4$\,s {\scriptsize(step 4{,}001)} & $6{,}309.1$\,s {\scriptsize(step 4{,}501)} & $6{,}981.7$\,s {\scriptsize(step 5{,}001)} & $\mathbf{3.4464}$ & $7{,}142.0$\,s \\
\bottomrule
\end{tabular}
}
\end{table}
